\par\Needspace{6\baselineskip}\medskip
\hypertarget{native-110-article081}{}\pdfbookmark[2]{81.}{article-081}
\begin{center}81.\end{center}
Si excentricitas $e$ data est, angulus $\Pi$ per aequationem
\[
\cos(A-\Pi)=-\frac{\cos(A-\tfrac12N-\tfrac12N')}{e\cos\tfrac12(N'-N)}
\]
invenietur, postquam angulus auxiliaris $A$ per aequationem
\[
\operatorname{tang}(A-\tfrac12N-\tfrac12N')=\frac{r'+r}{r'-r}\operatorname{tang}\tfrac12(N'-N)
\]
determinatus est. Ambiguitas in determinatione anguli $A-\Pi$ per ipsius cosinum remanens in natura problematis fundata est, ita ut problemati duabus solutionibus diversis satisfieri possit, e quibus quam adoptare quamve reiicere oporteat aliunde decidendum erit, ad quem finem valor saltem approximatus \SourceLeaf{111}ipsius $\Pi$ iam cognitus esse debet. — Postquam $\Pi$ inventus est, $p$ vel per formulas
\[
p=r(1+e\cos(N-\Pi))=r'(1+e\cos(N'-\Pi))
\]
vel per hanc computabitur
\[
p=\frac{2rr'e\sin\tfrac12(N'-N)\sin(\tfrac12N+\tfrac12N'-\Pi)}{r'-r}\mathpunct{\cdot}
\]
